\documentclass[12pt,a4paper]{article}

\usepackage[utf8]{inputenc}
\usepackage[T1]{fontenc}
\usepackage[polish]{babel}
\usepackage{lmodern}
\usepackage[margin=2.5cm]{geometry}
\usepackage{amsmath}
\usepackage{graphicx}
\usepackage{listings}
\usepackage{xcolor}
\usepackage{booktabs}
\usepackage{float}
\usepackage{hyperref}
\usepackage{caption}

\hypersetup{
    colorlinks=true,
    linkcolor=black,
    urlcolor=blue
}

\lstset{
    basicstyle=\ttfamily\small,
    keywordstyle=\color{blue},
    commentstyle=\color{gray},
    stringstyle=\color{red},
    breaklines=true,
    frame=single,
    numbers=left,
    numberstyle=\tiny\color{gray},
    showstringspaces=false
}

\title{Struktury Danych \\ Miniprojekt 2 \\ Kolejka priorytetowa}
\author{Paweł Bijak}
\date{\today}

\begin{document}

\maketitle

\section{Cel projektu}

Celem projektu jest zaimplementowanie własnymi środkami kolejki priorytetowej
typu MAX w języku C++ w oparciu o dwie różne struktury danych oraz przeprowadzenie
pomiarów czasu wykonania podstawowych operacji w celu praktycznej weryfikacji
ich teoretycznej złożoności obliczeniowej. Wybrane struktury to:
\begin{itemize}
    \item kopiec binarny przechowywany w tablicy,
    \item jednokierunkowa lista wiązana posortowana malejąco po priorytecie.
\end{itemize}

\section{Charakterystyka kolejki priorytetowej}

Kolejka priorytetowa to abstrakcyjny typ danych, w którym każdy element posiada
dodatkową wartość zwaną priorytetem. Kolejność obsługi elementów wyznacza
priorytet, a nie moment dodania do struktury. W wariancie MAX pierwszeństwo ma
element o najwyższym priorytecie.

W projekcie każdy element jest reprezentowany przez strukturę zawierającą:
\begin{itemize}
    \item wartość elementu (\texttt{int}),
    \item priorytet (\texttt{int}),
    \item numer wstawienia (\texttt{long long}) -- pole pomocnicze dla strategii FIFO.
\end{itemize}

\subsection{Strategia FIFO przy równych priorytetach}

W przypadku wielu elementów o tym samym priorytecie przyjęto strategię FIFO
(first in, first out). Każdy element przy wstawianiu otrzymuje rosnący numer
wstawienia. Przy porównaniu dwóch elementów o tym samym priorytecie
pierwszeństwo ma ten o mniejszym numerze wstawienia (czyli wstawiony wcześniej).

W przypadku operacji \texttt{modifyKey} przyjęto, że modyfikowany element
zachowuje swój pierwotny numer wstawienia. Oznacza to, że element po zmianie
priorytetu nie traktuje się jako ,,nowy''. Jest to spójne z intuicją, że
zmieniamy jedynie priorytet istniejącego elementu, a nie wstawiamy go na nowo.

\section{Implementacje}

\subsection{Kopiec binarny}

Kopiec binarny to drzewo binarne spełniające własność kopca: priorytet rodzica
jest większy lub równy priorytetowi każdego z jego dzieci (kopiec MAX).
W~projekcie kopiec został zaimplementowany jako tablica dynamiczna z indeksowaniem
0--bazowym, gdzie:
\begin{itemize}
    \item rodzic węzła $i$ znajduje się pod indeksem $\lfloor (i-1)/2 \rfloor$,
    \item lewe dziecko węzła $i$ znajduje się pod indeksem $2i+1$,
    \item prawe dziecko węzła $i$ znajduje się pod indeksem $2i+2$.
\end{itemize}

W razie braku miejsca tablica jest powiększana dwukrotnie. Operacje przywracające
własność kopca to \texttt{kopcujWGore} (po wstawieniu elementu) oraz
\texttt{kopcujWDol} (po wyciągnięciu maksimum lub modyfikacji w dół).

\subsubsection*{Złożoność operacji}
\begin{itemize}
    \item \texttt{insert} -- $O(\log n)$ (przesuwanie nowego elementu do właściwej pozycji),
    \item \texttt{extractMax} -- $O(\log n)$ (przywrócenie własności kopca po usunięciu korzenia),
    \item \texttt{peek} -- $O(1)$ (korzeń tablicy),
    \item \texttt{modifyKey} -- $O(n)$ (liniowe wyszukiwanie elementu po wartości) plus $O(\log n)$ na naprawę kopca,
    \item \texttt{rozmiar} -- $O(1)$.
\end{itemize}

\subsection{Lista wiązana}

Lista wiązana jednokierunkowa posortowana malejąco po priorytecie. Element
o~najwyższym priorytecie znajduje się zawsze na początku listy (głowa). Przy
wstawianiu element trafia za wszystkie istniejące elementy o tym samym lub
wyższym priorytecie, co realizuje strategię FIFO.

\subsubsection*{Złożoność operacji}
\begin{itemize}
    \item \texttt{insert} -- $O(n)$ (znalezienie odpowiedniej pozycji),
    \item \texttt{extractMax} -- $O(1)$ (zdjęcie głowy listy),
    \item \texttt{peek} -- $O(1)$ (zwrócenie głowy bez modyfikacji),
    \item \texttt{modifyKey} -- $O(n)$ (wyszukanie i ponowne wstawienie),
    \item \texttt{rozmiar} -- $O(1)$ (utrzymywany licznik).
\end{itemize}

\section{Metodyka pomiarów}

\subsection{Środowisko}

Pomiary wykonano na komputerze osobistym pod systemem Windows 10. Kompilacji
dokonano kompilatorem g++ (MSYS2/UCRT64, wersja 15.2.0) z flagami
\texttt{-O2 -std=c++17}. Czas mierzono z wykorzystaniem \texttt{std::chrono::steady\_clock}
z dokładnością do nanosekund.

\subsection{Sposób generowania danych}

Wartości elementów losowano z zakresu $[-1000, 1000]$. Priorytety losowano
z~zakresu $[0, 5n - 1]$, gdzie $n$ jest aktualnym rozmiarem kolejki. Zakres
priorytetów został wybrany jako pięciokrotnie większy od rozmiaru struktury,
zgodnie z zaleceniem treści projektu (,,kilkukrotnie większy''). Rozmiary
kolejek dobrane do badań to: $1\,000$, $2\,000$, $5\,000$, $10\,000$, $20\,000$,
$50\,000$, $100\,000$.

\subsection{Schemat pomiaru}

Dla każdego rozmiaru $n$ wykonano \textbf{5 niezależnych powtórzeń}
całego eksperymentu, a wyniki uśredniono. W każdym powtórzeniu:
\begin{itemize}
    \item operacja \texttt{insert} -- mierzono łączny czas wstawienia $n$ elementów do pustej kolejki, a~wynik podzielono przez $n$ (czas amortyzowany pojedynczego wstawienia),
    \item operacja \texttt{peek}, \texttt{modifyKey}, \texttt{extractMax} -- na pełnej kolejce o~rozmiarze $n$ wykonywano serię $K = \min(1000, n)$ wywołań i mierzono ich łączny czas, a~następnie dzielono przez $K$,
    \item operacja \texttt{rozmiar} -- ze względu na bardzo małą złożoność wykonywano $100\,000$ wywołań i wynik dzielono przez tę liczbę.
\end{itemize}

Zgodnie z założeniami, dla operacji \texttt{peek} i \texttt{rozmiar}, których
wynik nie jest wykorzystywany, dodano akumulator z modyfikatorem \texttt{volatile}
aby zapobiec optymalizacji usuwającej wywołanie przez kompilator.

\section{Wyniki pomiarów}

\textit{Wykresy i analiza zostaną dodane po zebraniu pełnego zestawu wyników.}

\section{Wnioski}

\textit{Wnioski zostaną sformułowane na podstawie pomiarów.}

\end{document}
